\documentclass[10pt,a4paper]{amsart}
\usepackage{amsmath,amssymb,amsthm,graphicx}
% For commutative diagrams we use Xy-pic
\usepackage[all]{xy}

% We use 2cell for 2-commutative diagrams.
\xyoption{2cell}
% Theorem environments.
%
\theoremstyle{plain}
\newtheorem{theorem}[subsection]{Theorem}
\newtheorem{proposition}[subsection]{Proposition}
\newtheorem{lemma}[subsection]{Lemma}

\theoremstyle{definition}
\newtheorem{definition}[subsection]{Definition}
\newtheorem{example}[subsection]{Example}
\newtheorem{exercise}[subsection]{Exercise}
\newtheorem{situation}[subsection]{Situation}

\theoremstyle{remark}
\newtheorem{remark}[subsection]{Remark}
\newtheorem{remarks}[subsection]{Remarks}

\numberwithin{equation}{subsection}

% Macros
%
\def\lim{\mathop{\mathrm{lim}}\nolimits}
\def\colim{\mathop{\mathrm{colim}}\nolimits}
\def\Spec{\mathop{\mathrm{Spec}}}
\def\Hom{\mathop{\mathrm{Hom}}\nolimits}
\def\Ext{\mathop{\mathrm{Ext}}\nolimits}
\def\SheafHom{\mathop{\mathcal{H}\!\mathit{om}}\nolimits}
\def\SheafExt{\mathop{\mathcal{E}\!\mathit{xt}}\nolimits}
\def\Sch{\mathit{Sch}}
\def\Mor{\mathop{\mathrm{Mor}}\nolimits}
\def\Ob{\mathop{\mathrm{Ob}}\nolimits}
\def\Sh{\mathop{\mathit{Sh}}\nolimits}
\def\NL{\mathop{N\!L}\nolimits}
\def\CH{\mathop{\mathrm{CH}}\nolimits}
\def\proetale{{pro\text{-}\acute{e}tale}}
\def\etale{{\acute{e}tale}}
\def\QCoh{\mathit{QCoh}}
\def\Ker{\mathop{\mathrm{Ker}}}
\def\Im{\mathop{\mathrm{Im}}}
\def\Coker{\mathop{\mathrm{Coker}}}
\def\Coim{\mathop{\mathrm{Coim}}}

% Boxtimes
%
\DeclareMathSymbol{\boxtimes}{\mathbin}{AMSa}{"02}

%
% Macros for moduli stacks/spaces
%
\def\QCohstack{\mathcal{QC}\!\mathit{oh}}
\def\Cohstack{\mathcal{C}\!\mathit{oh}}
\def\Spacesstack{\mathcal{S}\!\mathit{paces}}
\def\Quotfunctor{\mathrm{Quot}}
\def\Hilbfunctor{\mathrm{Hilb}}
\def\Curvesstack{\mathcal{C}\!\mathit{urves}}
\def\Polarizedstack{\mathcal{P}\!\mathit{olarized}}
\def\Complexesstack{\mathcal{C}\!\mathit{omplexes}}
% \Pic is the operator that assigns to X its picard group, usage \Pic(X)
% \Picardstack_{X/B} denotes the Picard stack of X over B
% \Picardfunctor_{X/B} denotes the Picard functor of X over B
\def\Pic{\mathop{\mathrm{Pic}}\nolimits}
\def\Picardstack{\mathcal{P}\!\mathit{ic}}
\def\Picardfunctor{\mathrm{Pic}}
\def\Deformationcategory{\mathcal{D}\!\mathit{ef}}
\newenvironment{reference}{}{}
\newenvironment{slogan}{}{}
\newenvironment{history}{}{}
\usepackage{cleveref}
\crefname{originaldefinition}{original definition}{original definitions}
\Crefname{originaldefinition}{Original definition}{Original definitions}
\makeatletter
\newlabel{definition-picard-group}{{0C7H}{source}{}{sourcepointer-1}{}}
\newlabel{definition-picard-group@cref}{{[originaldefinition][0][]0C7H}{[1][0][]source}{}{}{}}
\makeatother
\crefname{originaldefinition}{original definition}{original definitions}
\Crefname{originaldefinition}{Original definition}{Original definitions}
\makeatletter
\newlabel{definition-top-genus}{{0C79}{source}{}{sourcepointer-2}{}}
\newlabel{definition-top-genus@cref}{{[originaldefinition][0][]0C79}{[1][0][]source}{}{}{}}
\makeatother
\crefname{originallemma}{original lemma}{original lemmas}
\Crefname{originallemma}{Original lemma}{Original lemmas}
\makeatletter
\newlabel{lemma-bound-wm}{{0C9W}{source}{}{sourcepointer-3}{}}
\newlabel{lemma-bound-wm@cref}{{[originallemma][0][]0C9W}{[1][0][]source}{}{}{}}
\makeatother
\crefname{originallemma}{original lemma}{original lemmas}
\Crefname{originallemma}{Original lemma}{Original lemmas}
\makeatletter
\newlabel{lemma-contract}{{0C77}{source}{}{sourcepointer-4}{}}
\newlabel{lemma-contract@cref}{{[originallemma][0][]0C77}{[1][0][]source}{}{}{}}
\makeatother
\crefname{originallemma}{original lemma}{original lemmas}
\Crefname{originallemma}{Original lemma}{Original lemmas}
\makeatletter
\newlabel{lemma-contract-picard-group}{{0C7J}{source}{}{sourcepointer-5}{}}
\newlabel{lemma-contract-picard-group@cref}{{[originallemma][0][]0C7J}{[1][0][]source}{}{}{}}
\makeatother
\crefname{originallemma}{original lemma}{original lemmas}
\Crefname{originallemma}{Original lemma}{Original lemmas}
\makeatletter
\newlabel{lemma-genus-nonnegative}{{0C7C}{source}{}{sourcepointer-6}{}}
\newlabel{lemma-genus-nonnegative@cref}{{[originallemma][0][]0C7C}{[1][0][]source}{}{}{}}
\makeatother
\crefname{originallemma}{original lemma}{original lemmas}
\Crefname{originallemma}{Original lemma}{Original lemmas}
\makeatletter
\newlabel{lemma-picard-T-and-A}{{0CE7}{source}{}{sourcepointer-7}{}}
\newlabel{lemma-picard-T-and-A@cref}{{[originallemma][0][]0CE7}{[1][0][]source}{}{}{}}
\makeatother
\crefname{originallemma}{original lemma}{original lemmas}
\Crefname{originallemma}{Original lemma}{Original lemmas}
\makeatletter
\newlabel{lemma-recurring-symmetric-integer}{{0C6X}{source}{}{sourcepointer-8}{}}
\newlabel{lemma-recurring-symmetric-integer@cref}{{[originallemma][0][]0C6X}{[1][0][]source}{}{}{}}
\makeatother
\title{For a numerical type of genus at least two the l-torsion of its Picard group is bounded by the genus for large l}
\author{The Stacks Project Authors}
\date{stacks-project:models.tex:local-snapshot-20261012; GNU FDL 1.2 or later}
\begin{document}
\maketitle
\noindent\textit{Editorial note.} The counted unit is the Stacks Project Proposition proposition-bound-picard-group at source lines 4285-4298 together with its proof at source lines 4300-4319, a straight adjacent pairing in the chapter file. The chapter is the upstream file models.tex; the pool directory carries the longer name stacks-models and is not used as the chapter name anywhere in this card.\par\medskip
\noindent\textit{Editorial note.} This card has no context span: the statement and proof are self-contained apart from the eight pointers, and the notions they use are recorded in the standard prerequisites.\par\medskip
\noindent\textit{Editorial note.} The retained proof names eight results of the same chapter. All are declared as non-load-bearing navigation pointers with their stable Stacks tags looked up in the shared tag index; five of them are counted by cards of this lane and none of their intervals is reproduced here.\par\medskip
\noindent\textit{Editorial note.} The retained proof span starts at line 4300, the line of the opening marker read from the source, and contains exactly one proof environment.\par\medskip
\noindent\textit{Editorial note.} No citation occurs in the retained spans, so no bibliography entry is supplied.\par\medskip
\noindent\textit{Allowed standard prerequisites (editorial).}\begin{itemize}
\item Numerical types, their Picard group, the bound on the weighted intersection numbers and the contraction step. A numerical type is a tuple of an integer and integer families of multiplicities, intersection numbers, weights and genera subject to the five defining conditions of the chapter, and its Picard group is the cokernel of the matrix of divided intersection numbers; the genus of the type is the alternating sum determined by those data, and the topological genus is the expression in terms of the number of vertices and the positive intersection pairs. The chapter bounds the weighted intersection numbers of a minimal type of genus at least two, proves that a type of nonpositive genus has Picard group the integers, that a minimal type has genus at least the topological genus and that the torsion of a Picard group can be read off through the divided intersection matrix. The author's own definitions and lemmas of the chapter are the allowed inputs.. The counted proposition is the arithmetic estimate that makes the classification of the minimal types of a fixed genus finite, and it is the last ingredient of that classification in this section.
\end{itemize}\medskip
\ifdefined\hypertarget\hypertarget{sourcepointer-1}{}\fi
\noindent\textit{Source reference.} Original definition 0C7H: focused-seeds/AGP/sources/stacks-snapshot-20261011/tags.txt, source lines 15332--15332. Navigation pointer to the chapter's definition of the Picard group of a numerical type, named by the retained statement; the cited definition is an allowed input and its text is not reproduced here. The number is the stable Stacks tag of the target, looked up in the shared tag index.\par\medskip
\ifdefined\hypertarget\hypertarget{sourcepointer-2}{}\fi
\noindent\textit{Source reference.} Original definition 0C79: focused-seeds/AGP/sources/stacks-snapshot-20261011/tags.txt, source lines 15324--15324. Navigation pointer to the chapter's definition of the topological genus, used by the retained proof in its comparison of genera; the cited definition is an allowed input and its text is not reproduced here. The number is the stable Stacks tag of the target, looked up in the shared tag index.\par\medskip
\ifdefined\hypertarget\hypertarget{sourcepointer-3}{}\fi
\noindent\textit{Source reference.} Original lemma 0C9W: focused-seeds/AGP/sources/stacks-snapshot-20261011/tags.txt, source lines 15416--15416. Navigation pointer invoked by the retained proof for the bound on the weighted intersection numbers of a minimal type, which is what restricts the type to finitely many shapes; that lemma is counted by a separate card of this lane and its interval is not reproduced here. The number is the stable Stacks tag of the target, looked up in the shared tag index.\par\medskip
\ifdefined\hypertarget\hypertarget{sourcepointer-4}{}\fi
\noindent\textit{Source reference.} Original lemma 0C77: focused-seeds/AGP/sources/stacks-snapshot-20261011/tags.txt, source lines 15322--15322. Navigation pointer invoked by the retained proof for the effect of contracting a minus-one index on the genus; that lemma is counted by a separate card of this lane and its interval is not reproduced here. The number is the stable Stacks tag of the target, looked up in the shared tag index.\par\medskip
\ifdefined\hypertarget\hypertarget{sourcepointer-5}{}\fi
\noindent\textit{Source reference.} Original lemma 0C7J: focused-seeds/AGP/sources/stacks-snapshot-20261011/tags.txt, source lines 15334--15334. Navigation pointer invoked by the retained proof for the comparison of the Picard groups under a contraction, which is what lets the proof assume the type minimal; the cited statement is an allowed input of the same section and its proof is not reproduced here. The number is the stable Stacks tag of the target, looked up in the shared tag index.\par\medskip
\ifdefined\hypertarget\hypertarget{sourcepointer-6}{}\fi
\noindent\textit{Source reference.} Original lemma 0C7C: focused-seeds/AGP/sources/stacks-snapshot-20261011/tags.txt, source lines 15327--15327. Navigation pointer invoked by the retained proof for the fact that a minimal type of genus at least two has genus at least the topological genus; that lemma is counted by a separate card of this lane and its interval is not reproduced here. The number is the stable Stacks tag of the target, looked up in the shared tag index.\par\medskip
\ifdefined\hypertarget\hypertarget{sourcepointer-7}{}\fi
\noindent\textit{Source reference.} Original lemma 0CE7: focused-seeds/AGP/sources/stacks-snapshot-20261011/tags.txt, source lines 15567--15567. Navigation pointer invoked by the retained proof for the comparison of the Picard group of a type with the cokernel of its intersection matrix; that lemma is counted by a separate card of this lane and its interval is not reproduced here. The number is the stable Stacks tag of the target, looked up in the shared tag index.\par\medskip
\ifdefined\hypertarget\hypertarget{sourcepointer-8}{}\fi
\noindent\textit{Source reference.} Original lemma 0C6X: focused-seeds/AGP/sources/stacks-snapshot-20261011/tags.txt, source lines 15312--15312. Navigation pointer invoked by the retained proof for the bound on the order of the torsion of the cokernel of the intersection matrix; that lemma is counted by a separate card of this lane and its interval is not reproduced here. The number is the stable Stacks tag of the target, looked up in the shared tag index.\par\medskip
\begin{proposition}
\label{proposition-bound-picard-group}
Let $g \geq 2$. For every numerical type $T$ of genus $g$
and prime number $\ell > 768g$ we have
$$
\dim_{\mathbf{F}_\ell} \Pic(T)[\ell] \leq g
$$
where $\Pic(T)$ is as in Definition \ref{definition-picard-group}.
If $T$ is minimal, then we even have
$$
\dim_{\mathbf{F}_\ell} \Pic(T)[\ell] \leq g_{top} \leq g
$$
where $g_{top}$ as in Definition \ref{definition-top-genus}.
\end{proposition}

\begin{proof}
Say $T$ is given by $n, m_i, a_{ij}, w_i, g_i$.
If $T$ is not minimal, then there exists a $(-1)$-index.
After replacing $T$ by an equivalent type we may assume
$n$ is a $(-1)$-index. Applying Lemma \ref{lemma-contract-picard-group}
we find $\Pic(T) \subset \Pic(T')$ where $T'$
is a numerical type of genus $g$ (Lemma \ref{lemma-contract})
with $n - 1$ indices. Thus we conclude by induction on $n$
provided we prove the lemma for minimal numerical types.

\medskip\noindent
Assume that $T$ is a minimal numerical type of genus $\geq 2$.
Observe that $g_{top} \leq g$ by Lemma \ref{lemma-genus-nonnegative}.
If $A = (a_{ij})$ then since $\Pic(T) \subset \Coker(A)$
by Lemma \ref{lemma-picard-T-and-A}. Thus it suffices to prove
the lemma for $\Coker(A)$.
By Lemma \ref{lemma-bound-wm} we see that $m_i|a_{ij}| \leq 768g$ for
all $i, j$.
Hence the result by Lemma \ref{lemma-recurring-symmetric-integer}.
\end{proof}

\end{document}
